\chapter*{Introduction générale}
\addcontentsline{toc}{chapter}{Introduction générale}
\minitoc[n] % <-- [n] force l'affichage même pour un chapitre non numéroté


\section*{Introduction}

Mon travail de thèse s’inscrit dans le domaine des atomes froids. Un gaz d’atomes ultra-froids est constitué de quelques dizaines à plusieurs millions d’atomes maintenus à des températures extrêmement basses, de sorte que leurs propriétés quantiques deviennent dominantes. L’intérêt scientifique pour ces systèmes repose sur leur capacité à constituer des plateformes idéales pour explorer la physique quantique à N corps. Ces gaz sont ainsi utilisés dans des applications variées, comme la métrologie de précision, le développement de processeurs quantiques numériques, ou encore comme simulateurs quantiques de systèmes à N corps dont le comportement est difficile à calculer numériquement en raison de la complexité exponentielle avec le nombre de particules.  

\medskip

Parmi ces systèmes, les gaz unidimensionnels (1D) présentent un intérêt particulier. Ils permettent de comparer les observations expérimentales à des modèles théoriques et numériques précis, et servent de modèles pour comprendre des matériaux de faible dimension tels que les fils quantiques supraconducteurs, les chaînes de spins magnétiques, ou les réseaux 1D de jonctions Josephson. En confinant les atomes dans des potentiels très étroits transversalement, on gèle certains degrés de liberté et les atomes ne peuvent se déplacer que le long d’une direction, créant ainsi des gaz unidimensionnels de bosons, de fermions ou de mélanges de composants. Dans l’expérience réalisée au \gls{LCF}, des atomes de rubidium 87 sont piégés dans un potentiel magnétique produit par des fils déposés sur une puce atomique.  

\medskip

La réduction de la dimensionnalité modifie profondément les propriétés des gaz. Par exemple, les fluctuations quantiques sont amplifiées en 1D et empêchent la formation d’une condensation de Bose-Einstein à température nulle. Certains systèmes 1D sont intégrables, c’est-à-dire qu’ils possèdent autant de quantités conservées que de degrés de liberté. Cette intégrabilité entraîne des dynamiques particulières : contrairement aux systèmes ergodiques, ces gaz ne relaxent pas vers un état thermique classique, mais vers un état caractérisé par des fonctions conservées, appelées distributions de rapidités.  

\medskip

Ma thèse s’est concentrée sur l’étude d’un gaz 1D de bosons avec interactions de contact répulsives, un système intégrable étudié initialement par Lieb et Liniger. Les états propres de ce modèle, obtenus via l’ansatz de Bethe, sont décrits par des rapidités, qui peuvent être interprétées comme les vitesses de quasi-particules à vie infinie. Dans la limite thermodynamique, l’état relaxé du système est entièrement caractérisé par la distribution de rapidités.  

\medskip

Durant cette thèse, j’ai réalisé une partie des expériences avec Léa. Dubois, doctorant ayant travaillé précédemment sur le même dispositif. J’ai contribué à l’étude expérimentale de la distribution de rapidités, ainsi qu’à l’analyse numérique des résultats via des simulations basées sur la théorie \gls{GHD}. Mon apport a notamment porté sur l’étude des fluctuations dans le cadre du formalisme \gls{GGE} (Generalized Gibbs Ensemble) et sur la mise en place d’un potentiel dipolaire permettant de contrôler et moduler la distribution atomique.  

Le mémoire présente ainsi les résultats obtenus sur la caractérisation expérimentale et numérique des gaz de bosons 1D, à l’équilibre comme hors équilibre, via la mesure de la distribution de rapidités spatialement résolue. Il est structuré en sept chapitres regroupés en trois parties distinctes, détaillant à la fois les aspects expérimentaux, numériques et théoriques de cette étude.


\section*{Présentation des chapitres}

Le mémoire est structuré en sept chapitres, chacun abordant un aspect précis de l’étude des gaz de bosons 1D. Les chapitres sont conçus pour être indépendants à la lecture, bien que certaines références croisées permettent de suivre le fil complet de la thèse.

\paragraph{Chapitre 1 : Modèle de Lieb-Liniger et ansatz de Bethe}  
Ce chapitre introduit de manière pédagogique le modèle de Lieb-Liniger (\gls{LL}) et l’approche de l’ansatz de Bethe (\gls{BA}). Il commence par un rappel des notions de première et deuxième quantification pour expliquer la structure de l’Hamiltonien de \gls{LL}, en particulier l’interaction de contact. Les états de Bethe sont présentés comme des états propres de l’Hamiltonien, associés aux opérateurs de nombre de particules et de quantité de mouvement, et servent de point de départ pour l’étude des charges conservées, développée dans le chapitre suivant.

\paragraph{Chapitre 2 : Introduction au GGE et formalisme TBA}  
Le chapitre 2 présente le formalisme du Generalized Gibbs Ensemble (\gls{GGE}) et introduit la Thermodynamic Bethe Ansatz (\gls{TBA}) à travers les travaux de Yang et Yang. On y détaille notamment le calcul des moyennes de charges locales et le formalisme nécessaire pour relier l’état intégrable aux distributions de rapidités. Les notions de fluctuations sont évoquées, mais leur étude détaillée est laissée au chapitre 4.

\paragraph{Chapitre 3 : Analyse des fluctuations et échantillonnage du GGE}  
Ce chapitre explore les fluctuations des rapidités. Les simulations de Monte-Carlo permettent de vérifier si l’on peut échantillonner correctement le GGE. Le principe de fluctuation-réponse est utilisé pour tester la cohérence des formules obtenues, dans le régime où le formalisme TBA est valide.

\paragraph{Chapitre 4 : Dynamique des gaz 1D avec GHD}  
Le chapitre 4 étend l’étude à la dynamique hors équilibre à l’aide de la théorie Hydrodynamique Généralisée (GHD). On y montre comment le formalisme du GGE se traduit dans la dynamique, et que les distributions de rapidités locales restent des quantités conservées, contrairement à l’approche classique de Gibbs.

\paragraph{Chapitre 5 : Dispositif expérimental}  
Ce chapitre présente les aspects expérimentaux de la thèse : lasers, puce atomique, DMD, confinement 1D. Deux expériences principales sont décrites : la première étudie l’évolution d’un nuage initialement piégé dans un potentiel harmonique, dans le régime de quasi-condensation de Bose (\gls{qBEC}); la seconde utilise le DMD pour résoudre spatialement la distribution de rapidités et pour préparer des états hors équilibre.

\paragraph{Chapitre 6 : Analyse des déformations de profil et simulations GHD}  
Le chapitre 6 se concentre sur l’étude des déformations du profil de densité et des résultats numériques issus des simulations GHD, permettant de comparer les prédictions théoriques avec les mesures expérimentales.

\paragraph{Chapitre 7 : Potentiel dipolaire et perspectives expérimentales}  
Le dernier chapitre présente la réflexion sur l’ajout d’un potentiel dipolaire et les perspectives expérimentales associées. On y discute comment ce potentiel permet de préparer des distributions atomiques plus complexes et d’étudier de nouveaux régimes dynamiques.

\medskip

\noindent Bien que chaque chapitre puisse être lu indépendamment, certaines dépendances existent pour approfondir la compréhension globale :  
\begin{itemize}
    \item Chapitre 4 se base sur le Chapitre 2, qui lui-même repose sur le Chapitre 1.  
    \item Chapitre 6 fait référence au Chapitre 2.  
    \item Chapitre 3 utilise les résultats du Chapitre 2.  
    \item Chapitre 7 s’appuie sur le Chapitre 6, qui se réfère au Chapitre 2.
\end{itemize}


%
%{\color{blue}
%\section*{Introduction}
%
%Cette thèse s’inscrit dans le domaine des atomes ultra-froids, des systèmes où les propriétés quantiques deviennent dominantes grâce à des températures extrêmement basses. Ces gaz constituent des plateformes idéales pour explorer la physique quantique à N corps, avec des applications allant de la métrologie de précision au développement de processeurs quantiques, et servent également de simulateurs pour des systèmes complexes difficilement accessibles par calcul numérique.
%
%Parmi ces systèmes, les gaz unidimensionnels (1D) sont particulièrement intéressants. Confinés dans des potentiels étroits transversalement, les atomes ne peuvent se déplacer que le long d’une direction, ce qui permet d’étudier des gaz de bosons, de fermions ou de mélanges atomiques dans un cadre contrôlé. Dans notre expérience au Laboratoire Charles Fabry, des atomes de rubidium 87 sont piégés par des potentiels magnétiques générés par des fils déposés sur une puce atomique.
%
%La dimensionnalité réduite confère à ces gaz des propriétés uniques. Les fluctuations quantiques sont renforcées et empêchent la formation d’une condensation de Bose-Einstein en 1D. Certains systèmes unidimensionnels sont intégrables, possédant autant de quantités conservées que de degrés de liberté, ce qui entraîne des dynamiques hors du cadre thermique classique. L’état relaxé est alors décrit par la distribution de rapidités, fonction conservée caractérisant l’ensemble des quasi-particules.
%
%Cette thèse se concentre sur un gaz 1D de bosons avec interactions de contact répulsives, modèle intégrable étudié initialement par Lieb et Liniger. Les états propres, obtenus via l’ansatz de Bethe, sont décrits par des rapidités, interprétables comme des vitesses de quasi-particules à vie infinie. La limite thermodynamique du système est entièrement caractérisée par la distribution de rapidités.
%
%Au cours de ce travail, j’ai participé à la réalisation expérimentale avec L. Dubois, en mesurant la distribution de rapidités spatialement résolue pour des gaz inhomogènes. Mon apport principal concerne l’analyse numérique des données via la théorie hydrodynamique généralisée (GHD), l’étude des fluctuations dans le cadre du GGE, et la mise en place d’un potentiel dipolaire permettant de moduler et contrôler la distribution atomique.
%
%Le mémoire présente les résultats obtenus sur la caractérisation expérimentale et numérique des gaz de bosons 1D, à l’équilibre comme hors équilibre, en mettant l’accent sur la distribution de rapidités localement résolue. Les sept chapitres sont organisés en trois parties distinctes, couvrant les aspects expérimentaux, numériques et théoriques de cette étude.
%}


